\section{Introduction}

A Sense/Compute/Control (SCC) application is one that interacts with
the environment~\cite{Tayl09a}. The SCC architectural pattern
guides the description of SCC applications and involves four kinds of
components, organized into layers~\cite{Cass11a,Edwar09a}: (1)
\emph{sensors} at the bottom, which obtain information about the
environment; (2) then \emph{context operators}, which process this
information; (3) then \emph{control operators}, which use this refined
information to control (4) \emph{actuators} at the top, which finally
impact the environment. A robotics application is a kind of SCC
application where the environment is composed of a robot
(sensors/actuators/body, control architecture, \etc{}) and the robot's
neighborhood (the walls, ground, people, \etc{})~\cite{Sicil08a}. As
noticed by Taylor \etal{}~\cite{Tayl09a}, the Sense/Plan/Act
architecture~\cite{Sicil08a}, widely used in robotics, closely
resembles the SCC architectural pattern.

\diaspec{} is a design language dedicated to describing SCC
applications. From such a design description, the \diaspec{} compiler
produces a dedicated Java programming framework
(Figure~\ref{fig:diaspec-process}) that is both \emph{prescriptive}
and \emph{restrictive}: it is prescriptive in the sense that it guides
the developer, and it is restrictive in the sense that it limits the
developer to what the design description allows. For each component
description in the design, the compiler generates an abstract class.
An abstract method in this class (\emph{hole} in
Figure~\ref{fig:diaspec-process}) represents code to be provided by
the developer to allow him to implement the application logic
(\textit{bump} in Figure~\ref{fig:diaspec-process}). To implement a
component, the developer implements these methods in a subclass of
this abstract class and thus never modifies generated code.

\begin{figure}
  \centering
  \includegraphics[width=\linewidth]{diaspec-process}
  \caption{Overview of the \diaspec{} development process}
  \label{fig:diaspec-process}
\end{figure}


By separating application logic (implemented by the developers) and
runtime support (generated in the programming framework), \diaspec{}
facilitates the design, implementation and evolution of SCC
applications. 

\subsection*{Contributions}

Our contributions are as follows:

\begin{itemize}
\item \emph{A report} on an experiment of designing and implementing a
  standard robotics application in the SCC architectural style with
  the \diaspec{} language (Sections~\ref{sec:designing}
  and~\ref{sec:implementing}). This report includes detailed
  instructions and guidelines to allow further experiments.
\item \emph{A discussion} of the benefits and problems of using \diaspec{} in
  a robotics setting (Section~\ref{sec:discussing}). This discussion
  includes a list of changes to \diaspec{} that would make it a better
  framework for developing new robotics applications.
\end{itemize}

We finally highlight some related works and conclude in
sections~\ref{sec:related} and~\ref{sec:conclusion}.

%%% Local Variables:
%%% mode: latex
%%% coding: utf-8
%%% TeX-master: "dslrob"
%%% TeX-PDF-mode: t
%%% ispell-local-dictionary: "english"
%%% End:

% LocalWords:  SCC
